\addcontentsline{toc}{section}{Notation}
\section*{Notation}
\subsection*{Allgemein}
\begin{tabbing}
xxxxxxxxxxxxxxxxxx\=xxxxxxxxxxxxxxxxxxxxxxxxxxxxxxxxxxxxx\kill
$t$, $u$, $x$, $z$, $\theta$, $\theta_0$, $\theta_1$ \>Reelle Parameter \\
$i, j, k, l, m, n$ \>Natürliche Zahlen \\
$\arctan$ \>Arkustangens \\
$\erf$ \>Gaußsche Fehlerfunktion \\
$g(x)$, $\phi(x)$ \>Funktionen \\
\end{tabbing}
\subsection*{Graphen}
\begin{tabbing}
xxxxxxxxxxxxxxxxxx\=xxxxxxxxxxxxxxxxxxxxxxxxxxxxxxxxxxxxx\kill
$b(G)$ \>Brennzahl von $G$ \\
$b_i$ \>$i$-ter Knoten einer Brennsequenz \\
$b_{i,j}, s_{i,j}$, $x_j$ \>Binärvariablen der ILP-PROP-Formulierung \\
$D$ \>Durchmesser des Graphen $G$ \\
$E$ \>Kantenmenge \\
$G = (V, E)$ \>Graph \\
$H_s(v)$ \>Höhe des Knotens $v$ im BFS-Baum mit Startknoten $s$ \\
$K_n$ \>Vollständiger Graph mit $n$ Knoten \\
$m$ \>Kantenanzahl \\
$n$ \>Knotenanzahl \\
$N(v)$, $N_i(v)$ \>1-Umgebung bzw. $i$-Umgebung des Knotens $v$ \\
$s, u, u_i, v, v_i$ \>Knoten \\
$P_n$ \>Pfadgraph mit $n$ Knoten \\
$U$ \>Obere Schranke für die Brennzahl in der ILP-PROP-Formulierung\\
$V$ \>Knotenmenge \\
$w$ \>Kantengewicht \\
$Z(v)$ \>Zusammenhangskomponente des Knotens $v$ \\
\end{tabbing}
\subsection*{Zufallsvariablen, Verteilungen und Statistik}
\begin{tabbing}
xxxxxxxxxxxxxxxxxx\=xxxxxxxxxxxxxxxxxxxxxxxxxxxxxxxxxxxxx\kill
$X$, $Y$, $X_i$, $Y_i$ \>Zufallsvariablen \\
$\alpha$ \>Signifikanzniveau \\
$t_{\mathdutch{p}, \mathdutch{n}}$ \>$\mathdutch{p}$-Quantil der $t$-Verteilung mit $\mathdutch{n}$ Freiheitsgraden \\
$C$ \>Berry-Esseen-Konstante \\
$\chi^2_{\mathdutch{p}, n}$ \>$\mathdutch{p}$-Quantil der $\chi^2$-Verteilung mit $\mathdutch{n}$ Freiheitsgraden \\
$\E{X}$ \>Erwartungswert von $X$ \\
$\mu$, $\sigma$ \>Parameter der Normalverteilung \\
$\mu_X$, $\sigma_X$ \>Parameter der Normalverteilung von $X$ \\
$\gamma$ \>Schiefe \\
$\epsilon$ \>Fehlerterm \\
$f$ \>\pdf einer Normalverteilung \\
$F$ \>Verteilungsfunktion einer Normalverteilung \\
$f_X$ \>\pdf von $X$ \\
$F_X$ \>\cdf von $X$ \\
\end{tabbing}

